\documentclass{ximera}



\title{Fractions:  Adding and Subtracting}   % Use xmlatex all -s to compile when SERVE refuses to work.
\author{Kelly Stady}
\license{CC: 0}

\begin{document}


\begin{abstract}

\end{abstract}

\maketitle

label{xim:FractionsAddSub}


\section*{A ``Common" Mistake}

A student is asked to add $\textstyle \frac{1}{2}$ and $\textstyle \frac{1}{3}$ and writes the following.  Is their work correct?


\[ \frac{1}{2} + \frac{1}{3} = \frac{2}{5} \]


To answer this question, let’s look at a visual representation of the student's logic.


\begin{center}
\begin{tikzpicture}[scale=0.8, line width=0.8pt, font=\Large]

% --- 1/2 rectangle ---
\draw (0,0) rectangle (2,2);
\draw (1,0) -- (1,2);
\fill[glaucous,opacity=0.5] (0,0) rectangle (1,2);
\node at (1,-0.8) {$\mathbf{\frac{1}{2}}$};

% plus sign
\node at (3,-0.8) {$+$};

% --- 1/3 rectangle ---
\draw (4,0) rectangle (6,2);
\draw (4+2/3,0) -- (4+2/3,2);
\draw (4+4/3,0) -- (4+4/3,2);
\fill[orange!80!black,opacity=0.5] (4,0) rectangle (4+2/3,2);
    \node at (5,-0.8) {$\mathbf{\frac{1}{3}$};

% not equal
\node at (7,-0.8) {$\overset{\text{\scalebox{2.5}{?}}}{=}$};

%\node at (7,-0.8) {$\neq$};

% --- 2/5 rectangle ---
\draw (8,0) rectangle (10,2);
\foreach \x in {1,2,3,4} {
  \draw (8+2*\x/5,0) -- (8+2*\x/5,2);
}
\fill[teal,opacity=0.5] (8,0) rectangle (8+4/5,2);
\node at (9,-0.8) {$\mathbf{\frac{2}{5}}$};

\end{tikzpicture}
\xmalt{Three squares are shown. The first is divided into 2 equal vertical parts with 1 shaded, labeled one-half. A plus sign appears. The second is divided into 3 equal vertical parts with 1 shaded, labeled one-third. An equal sign with a question mark above it appears. The third square is divided into 5 equal vertical parts with 2 shaded, labeled two-fifths.}
\end{center}


Since $\textstyle \frac{1}{3}$ is \emph{added} to $\textstyle \frac{1}{2},$ the result should be greater than $\textstyle \frac{1}{2}.$  
However, the student's answer of $\textstyle \frac{2}{5}$ is actually less than $\textstyle \frac{1}{2}.$  This shows that simply 
adding the numerators and denominators doesn't work!  To get the correct answer, we need to divide each whole into the same number of
pieces. 


\begin{center}
\textbf{Fractions can only be added or subtracted when the pieces are the same size.}
\end{center}

Recall that the denominator of a fraction tells you how many equal-sized pieces a whole has been divided into.  If we decide to divide each whole into six equal pieces, the original fractions can be rewritten as follows.

\begin{align*}
\frac{1}{2} &= \frac{1}{2} \cdot \frac{3}{3} & \frac{1}{3} &= \frac{1}{3} \cdot \frac{2}{2} \\[10pt]
            &= \frac{3}{6}                   &             &= \frac{2}{6}
\end{align*}

Now we can visually represent the sum $\textstyle \frac{1}{2} + \frac{1}{3}$ as an equivalent sum with equal pieces:

\begin{center}
\begin{tikzpicture}[scale=0.8, line width=0.8pt, font=\Large]

% --- 3/6 rectangle (1/2 equivalent) ---
\draw (0,0) rectangle (2,2);
\foreach \x in {1,2,3,4,5} {
  \draw (0+2*\x/6,0) -- (0+2*\x/6,2);
}
\fill[glaucous,opacity=0.5] (0,0) rectangle (1,2);
\node at (1,-0.8) {$\mathbf{\frac{3}{6}}$};

% plus
\node at (3,-0.8) {$+$};

% --- 2/6 rectangle (1/3 equivalent) ---
\draw (4,0) rectangle (6,2);
\foreach \x in {1,2,3,4,5} {
  \draw (4+2*\x/6,0) -- (4+2*\x/6,2);
}
\fill[orange!80!black,opacity=0.5] (4,0) rectangle ({4 + 2/3},2);
\node at (5,-0.8) {$\mathbf{\frac{2}{6}}$};

% equals
\node at (7,-0.8) {$=$};

% --- 5/6 result ---
\draw (8,0) rectangle (10,2);
\foreach \x in {1,2,3,4,5} {
  \draw (8+2*\x/6,0) -- (8+2*\x/6,2);
}
\fill[teal, opacity=0.5] (8,0) rectangle ({8+5/3},2);
\node at (9,-0.8) {$\mathbf{\frac{5}{6}}$};

\end{tikzpicture}
\xmalt{Three rectangles are shown, each divided into 6 equal vertical parts. The first has 3 parts shaded, labeled three-sixths. A plus sign appears. The second has 2 parts shaded, labeled two-sixths. An equals sign appears. The third has 5 parts shaded, labeled five-sixths, showing the correct sum.}
\end{center}


Combining 3 sixths and 2 sixths results in a total of 5 sixths.  The entire problem is written mathematically below:

\begin{align*}
\frac{1}{2} + \frac{1}{3} &= \frac{1}{2}\cdot \frac{3}{3} + \frac{1}{3} \cdot \frac{2}{2} \\[10pt]
&= \frac{3}{6} + \frac{2}{6} \\[10pt]
&= \frac{5}{6}
\end{align*}



\begin{remark}
To learn mathematics, your solutions should be written using clear, logical steps;  this process will deepen your understanding, prevent errors and serve as a test of your knowledge---if you cannot clearly communicate your work, you have not yet fully grasped the concept.
\end{remark}

\begin{procedure}[\textbf{Adding and Subtracting Fractions}]

~ \\[1em]

\begin{description}
  \item[Find a Common Denominator:] Before adding or subtracting, the pieces must be the same size. If the denominators are different, rewrite the fractions as equivalent fractions with a common denominator.

  \item[Add or Subtract the Numerators:] This tells you how many pieces you have.

  \item[Keep the Denominator the Same:] The size of the pieces doesn't change during addition or subtraction, so the denominator stays the same.

  \item[Simplify:]  Check to see if the resulting fraction can be simplified.  If so, simplify.
\end{description}

\end{procedure}


\begin{problem}
Use the visual representation to find the difference.


%\[ \frac{4}{5} - \frac{1}{5} = \frac{\answer{3}}{\answer{5}} \]

\begin{center}
\begin{tikzpicture}[scale=0.8, line width=0.8pt]
% --- 4/5 rectangle ---
\draw (0,0) rectangle (2,2);
\foreach \x in {1,2,3,4} { \draw (2*\x/5,0) -- (2*\x/5,2); }
\fill[glaucous,opacity=0.5] (0,0) rectangle (1.6,2); % 4/5 of 2 units is 1.6
\node at (1,-0.8) {$\frac{4}{5}$};

% minus sign (at label level)
\node at (3,-0.8) {$-$};

% --- 1/5 rectangle ---
\draw (4,0) rectangle (6,2);
\foreach \x in {1,2,3,4} { \draw (4+2*\x/5,0) -- (4+2*\x/5,2); }
\fill[orange!80!black,opacity=0.5] (4,0) rectangle (4.4,2); % 1/5 of 2 units is 0.4
\node at (5,-0.8) {$\frac{1}{5}$};

% equals sign (at label level)
\node at (7,-0.8) {$=$};

% --- Result rectangle (3/5) ---
\draw (8,0) rectangle (10,2);
\foreach \x in {1,2,3,4} { \draw (8+2*\x/5,0) -- (8+2*\x/5,2); }
\fill[teal,opacity=0.5] (8,0) rectangle (9.2,2); % 3/5 of 2 units is 1.2

% Giant question mark
\node at (9,-0.8) {\textup{\scalebox{2}{?}}};

\end{tikzpicture}
\end{center}

$\frac{4}{5} - \frac{1}{5} = \frac{\answer{3}}{\answer{5}}$

\end{problem}


\section*{The Least Common Denominator}

When adding or subtracting fractions with different denominators, we will use the \emph{smallest} or \emph{least common denominator}.  This will
make the simplification at the end as easy as possible.

\begin{definition}[Least Common Denominator (LCD)]

The \textbf{least common denominator} is the smallest number that is a multiple of two or more denominators.  Equivalently, it is the smallest number that is
evenly divisible by two or more denominators.

\end{definition}

\begin{remark}
The \textbf{multiples of a number} are found by multiplying the number by the counting numbers: $\{1, 2, 3, 4, \ldots \}.$  For example, the first three multiples of 4 are:
4, 8 and 12.

\begin{problem}
List the first five multiples of $3$, separated by commas.
\[
\answer[format=set]{3, 6, 9, 12, 15}
\]
\end{problem}

\end{remark}


\begin{procedure}[Finding the Least Common Denominator]


Write the multiples  of each denominator until you find the smallest number that appears in each list.  \textbf{Alternatively}, you can list
the multiples of only the largest denominator and stop when you find the smallest number that is evenly divisible by the other denominators.
\end{procedure}


\begin{example}
Find the LCD of $\textstyle \frac{3}{4}$ and $\textstyle \frac{1}{7}$ using the list of multiples given below.


\begin{align*}
\textbf{Multiples of 4: } & 4, 8, 12, 16, 20, 24, \mathbf{28}, 32, 36, \ldots \\
\textbf{Multiples of 7: } & 7, 14, 21, \mathbf{28}, 35, 42, 49, 56, 63, \ldots
\end{align*}

The smallest number that appears in both lists is 28, so the LCD = 28.  Note that we could have saved ourselves some time by
only listing the multiples of 7 and looking for the smallest number that is evenly divisible by 4.

\begin{problem}
Now that you know the LCD = 28.  Rewrite each fraction as an equivalent fraction with a denominator of 28.

\[ \frac{3}{4} = \frac{\answer{21}}{\answer{28}} \]

\[ \frac{1}{7} = \frac{\answer{4}}{\answer{28}} \]

\end{problem}


\end{example}


\begin{problem}
A community land trust is dividing a large lot to combat food insecurity. They decide to use $\textstyle \frac{2}{3}$ of the lot for growing fresh vegetables to donate to local food pantries and $\textstyle \frac{1}{5}$ of the lot for a fruit orchard. What total fraction of the land is dedicated to food production?

    First, find the Least Common Denominator (LCD) for $\frac{2}{3}$ and $\frac{1}{5}.$

    \begin{hint}
      To find the LCD, look for the smallest number that is a multiple of both 3 and 5:
      \begin{align*}
          \textbf{Multiples of 3: } & 3, 6, 9, 12, 15, 18, 21, \dots \\
          \textbf{Multiples of 5: } & 5, 10, 15, 20, 25, 30, \dots
      \end{align*}
    \end{hint}

    The LCD is $\answer{15}$.

    Now, use the LCD to write each fraction as an equivalent fraction:
    \begin{align*}
        \text{Vegetable Garden: } \frac{2}{3} &= \frac{\answer{10}}{\answer{15}} \\
                                                & \\
        \text{Fruit Orchard: } \frac{1}{5} &= \frac{\answer{3}}{\answer{15}}
    \end{align*}

    Finally, add the equivalent fractions.
    \[
     \textrm{Fraction of land used for food} \ = \answer{\frac{13}{15}}
    \]

\end{problem}



\begin{problem}
Factory farms store animal waste in massive open-air pits called lagoons. In one community, a lagoon is $\textstyle \frac{5}{6}$ full of untreated waste. To lower the level and prevent an overflow into the local groundwater, the farm operators spray an amount equal to $\textstyle \frac{1}{4}$ of the lagoon’s total capacity into the air and onto nearby fields.  What fraction of the lagoon's total capacity remains full?

First, find the Least Common Denominator (LCD) for $\frac{5}{6}$ and $\frac{1}{4}$.
\begin{hint}
To find the LCD, write the multiples of 6 and look for the smallest number that is evenly divisible by 4. Or write the multiples of both numbers and look for the smallest number that appears in both lists.
\end{hint}

The LCD is $\answer{12}$.

Now, use the LCD to write each fraction as an equivalent fraction:
\begin{align*}
\text{Initial Waste level: } \frac{5}{6} &= \frac{\answer{10}}{\answer{12}} \\
& \\
\text{Waste removed: } \frac{1}{4} &= \frac{\answer{3}}{\answer{12}}
\end{align*}

Finally, subtract the equivalent fractions.
\[
\textrm{Remaining waste level: }  = \frac{ \answer{7} }{ \answer{12} }
\]

\begin{feedback}
While spraying the waste on fields prevents the lagoon from overflowing, it creates a new problem for the community. The waste is often sprayed as a fine mist that can travel for miles, carrying harmful bacteria, nitrates and intense odors into nearby neighborhoods.
For the people living nearby, this isn't just an unpleasant smell---it is a serious health hazard! Studies have shown that families living near these spray fields experience higher rates of asthma, respiratory infections and even contaminated well water. This is a clear example of environmental injustice, as these industrial farms are often located near low-income communities who have the fewest resources to fight back against the pollution.
\end{feedback}
\end{problem}



\begin{question}
    A city-wide study on workforce equity found that $\textstyle \frac{1}{6}$ of the city's engineering jobs were held by women. After a local mentoring program was launched, the fraction of engineering jobs held by women increased by $\textstyle \frac{1}{10}$. What fraction of engineering jobs are now held by women?

    First, find the least common denominator (LCD)for $\frac{1}{6}$ and $\frac{1}{10}$.

    %\begin{hint}
      %To find the LCD, look for the smallest number that is a multiple of both 6 and 10:
      %\begin{align*}
         % \textbf{Multiples of 6: } & 6, 12, 18, 24, \mathbf{30}, 36, \dots \\
          %\textbf{Multiples of 10: } & 10, 20, \mathbf{30}, 40, \dots
      %\end{align*}
    %\end{hint}

    The LCD is $\answer{30}$.

    Now, use the LCD to write each fraction as an equivalent fraction:
    \begin{align*}
        \text{Original jobs: } \frac{1}{6} &= \frac{\answer{5}}{\answer{30}} \\
        & \\
        \text{Increase: } \frac{1}{10} &= \frac{\answer{3}}{\answer{30}}
    \end{align*}

    Add the equivalent fractions to find the total:
    \[
      \textrm{Fraction of engineering jobs held by women} = \frac{\answer{8}}{\answer{30}}
    \]

    Finally, simplify your answer:
    \[
      \textrm{Fraction of engineering jobs held by women} = \frac{\answer{4}}{\answer{15}}
    \]
\end{question}


\end{document} 